Artificial intelligence has become a prominent theme in corporate disclosure, yet it is difficult for investors to distinguish statements that reflect genuine implementation from those that primarily amount to narrative positioning. We classify AI-related sentences in annual 10-K filings as actionable, speculative, or irrelevant, validate this disclosure layer against subsequent AI patenting, and combine disclosure composition with contemporaneous patent weakness to construct PatentMismatch, a scalable, audited measure of AI-washing. In our 2016--2025 sample of U.S. public firms, PatentMismatch rises sharply in the post-2022 period and characterizes 41.5\% of AI-talking firm-years in 2024 and 42.1\% in 2025. Markets do not appear to fully discipline this at the filing date. Instead, the stronger correction emerges after disclosure, where a long--short strategy that buys firms with more credible AI disclosure and shorts mismatch firms earns its strongest return over the subsequent quarter, with an equal-weight alpha of 0.93\% per month, or about 11.2\% annualized. The cross-sectional valuation effects are stronger outside the largest firms. Actionable AI language aligns more closely with subsequent AI patenting than speculative language, and low-credibility disclosure predicts weaker future AI patent output. The evidence suggests that AI-washing is economically relevant both because it departs from later technological realization and because markets do not appear to price that gap immediately.